\documentclass[12pt,a4paper]{article}
\usepackage[UTF8]{ctex}
\usepackage{amsmath,amssymb,amsthm}
\usepackage{geometry}

\geometry{margin=2.5cm}

\title{约束力$F_{\text{constraint}} = A^T(\theta)\lambda$的详细解释}
\author{第11章补充说明}
\date{}

\begin{document}

\maketitle

\section{问题提出}

在约束动力学中，约束力表示为：
\begin{equation}
F_{\text{constraint}} = A^T(\theta)\lambda
\end{equation}

\textbf{问题}：
\begin{enumerate}
\item 为什么约束力是这种形式？
\item $F_{\text{constraint}}$的矩阵构型是什么？
\item 为什么是$A^T$而不是$A$？
\item 拉格朗日乘数$\lambda$的物理意义是什么？
\end{enumerate}

\section{矩阵构型分析}

\subsection{维度分析}

\textbf{已知}：
\begin{itemize}
\item $A(\theta) \in \mathbb{R}^{k \times 6}$：约束矩阵（$k$个约束，6维任务空间）
\item $\lambda \in \mathbb{R}^k$：拉格朗日乘数向量（$k$个约束对应$k$个乘数）
\item $F_{\text{constraint}} \in \mathbb{R}^6$：约束力（6维wrench）
\end{itemize}

\textbf{矩阵乘法}：
\begin{equation}
F_{\text{constraint}} = A^T(\theta)\lambda
\end{equation}

\textbf{维度检查}：
\begin{align}
A^T(\theta) &\in \mathbb{R}^{6 \times k} \quad \text{（$A$的转置）} \\
\lambda &\in \mathbb{R}^k \\
F_{\text{constraint}} &= A^T(\theta)\lambda \in \mathbb{R}^6
\end{align}

\textbf{验证}：$(6 \times k) \times (k \times 1) = 6 \times 1$ ✓

\subsection{矩阵结构}

\textbf{约束矩阵$A(\theta)$的结构}：
\begin{equation}
A(\theta) = \begin{bmatrix}
a_{11} & a_{12} & \cdots & a_{16} \\
a_{21} & a_{22} & \cdots & a_{26} \\
\vdots & \vdots & \ddots & \vdots \\
a_{k1} & a_{k2} & \cdots & a_{k6}
\end{bmatrix} \in \mathbb{R}^{k \times 6}
\end{equation}

\textbf{转置矩阵$A^T(\theta)$的结构}：
\begin{equation}
A^T(\theta) = \begin{bmatrix}
a_{11} & a_{21} & \cdots & a_{k1} \\
a_{12} & a_{22} & \cdots & a_{k2} \\
\vdots & \vdots & \ddots & \vdots \\
a_{16} & a_{26} & \cdots & a_{k6}
\end{bmatrix} \in \mathbb{R}^{6 \times k}
\end{equation}

\textbf{拉格朗日乘数向量$\lambda$}：
\begin{equation}
\lambda = \begin{bmatrix}
\lambda_1 \\
\lambda_2 \\
\vdots \\
\lambda_k
\end{bmatrix} \in \mathbb{R}^k
\end{equation}

\textbf{约束力$F_{\text{constraint}}$}：
\begin{equation}
F_{\text{constraint}} = A^T(\theta)\lambda = \begin{bmatrix}
\sum_{i=1}^k a_{i1}\lambda_i \\
\sum_{i=1}^k a_{i2}\lambda_i \\
\vdots \\
\sum_{i=1}^k a_{i6}\lambda_i
\end{bmatrix} \in \mathbb{R}^6
\end{equation}

\section{为什么是$A^T$而不是$A$？}

\subsection{约束子空间和力子空间}

\textbf{约束子空间}：
\begin{itemize}
\item 约束$A(\theta)V = 0$定义了速度$V$必须位于的约束子空间
\item 约束子空间的维度：$6 - k$（6维空间减去$k$个约束）
\item 约束子空间由$A(\theta)$的零空间（null space）定义
\end{itemize}

\textbf{力子空间}：
\begin{itemize}
\item 约束力$F_{\text{constraint}}$必须位于约束子空间的正交补空间
\item 正交补空间的维度：$k$
\item 正交补空间由$A^T(\theta)$的列空间（column space）定义
\end{itemize}

\subsection{数学原理}

\textbf{基本事实}：
\begin{itemize}
\item 如果$A(\theta)V = 0$（约束子空间）
\item 那么$F_{\text{constraint}}$必须与约束子空间正交
\item 即$F_{\text{constraint}}^T V = 0$对所有满足$A(\theta)V = 0$的$V$成立
\end{itemize}

\textbf{推导}：

如果$A(\theta)V = 0$，那么$V$在$A(\theta)$的零空间中。

约束力$F_{\text{constraint}}$必须与所有满足约束的速度$V$正交：
\begin{equation}
F_{\text{constraint}}^T V = 0 \quad \text{对所有满足$A(\theta)V = 0$的$V$}
\end{equation}

这意味着$F_{\text{constraint}}$在$A(\theta)$的行空间（row space）中，即$A^T(\theta)$的列空间中。

因此：
\begin{equation}
F_{\text{constraint}} = A^T(\theta)\lambda
\end{equation}

其中$\lambda$是拉格朗日乘数，表示约束力在各个约束方向上的分量。

\subsection{几何解释}

\textbf{6维力空间的分解}：

\begin{itemize}
\item \textbf{约束子空间}（维度$6-k$）：
\begin{itemize}
\item 由$A(\theta)$的零空间定义
\item 速度$V$可以自由运动的方向
\item 约束力在这个子空间中为零
\end{itemize}

\item \textbf{约束力的子空间}（维度$k$）：
\begin{itemize}
\item 由$A^T(\theta)$的列空间定义
\item 约束力$F_{\text{constraint}}$必须位于这个子空间
\item 这是约束子空间的正交补空间
\end{itemize}
\end{itemize}

\textbf{正交性}：
\begin{itemize}
\item 约束子空间和约束力子空间是正交的
\item 因此约束力不会影响自由运动方向
\item 约束力只影响约束方向
\end{itemize}

\section{拉格朗日乘数$\lambda$的物理意义}

\subsection{基本定义}

\textbf{拉格朗日乘数$\lambda$}：
\begin{itemize}
\item $\lambda \in \mathbb{R}^k$：$k$个标量，每个对应一个约束
\item $\lambda_i$：第$i$个约束的拉格朗日乘数
\item 表示第$i$个约束产生的约束力的大小
\end{itemize}

\subsection{物理意义}

\textbf{约束力的分解}：
\begin{equation}
F_{\text{constraint}} = A^T(\theta)\lambda = \sum_{i=1}^k \lambda_i \cdot \text{第$i$列}(A^T)
\end{equation}

\textbf{解释}：
\begin{itemize}
\item $A^T(\theta)$的每一列表示一个约束方向
\item $\lambda_i$表示第$i$个约束方向上的约束力大小
\item 约束力$F_{\text{constraint}}$是各个约束方向上的约束力的线性组合
\end{itemize}

\subsection{具体例子}

\textbf{例子：擦黑板（$k = 3$个约束）}

\textbf{约束矩阵}：
\begin{equation}
A(\theta) = \begin{bmatrix}
1 & 0 & 0 & 0 & 0 & 0 \\  % $\omega_x = 0$
0 & 1 & 0 & 0 & 0 & 0 \\  % $\omega_y = 0$
0 & 0 & 0 & 0 & 0 & 1     % $v_z = 0$
\end{bmatrix}
\end{equation}

\textbf{转置矩阵}：
\begin{equation}
A^T(\theta) = \begin{bmatrix}
1 & 0 & 0 \\
0 & 1 & 0 \\
0 & 0 & 0 \\
0 & 0 & 0 \\
0 & 0 & 0 \\
0 & 0 & 1
\end{bmatrix}
\end{equation}

\textbf{拉格朗日乘数}：
\begin{equation}
\lambda = \begin{bmatrix}
\lambda_1 \\  % 对应$\omega_x = 0$的约束
\lambda_2 \\  % 对应$\omega_y = 0$的约束
\lambda_3     % 对应$v_z = 0$的约束
\end{bmatrix}
\end{equation}

\textbf{约束力}：
\begin{align}
F_{\text{constraint}} &= A^T(\theta)\lambda \\
&= \begin{bmatrix}
1 & 0 & 0 \\
0 & 1 & 0 \\
0 & 0 & 0 \\
0 & 0 & 0 \\
0 & 0 & 0 \\
0 & 0 & 1
\end{bmatrix} \begin{bmatrix}
\lambda_1 \\
\lambda_2 \\
\lambda_3
\end{bmatrix} \\
&= \begin{bmatrix}
\lambda_1 \\  % $m_x$（绕$x$轴的力矩）
\lambda_2 \\  % $m_y$（绕$y$轴的力矩）
0 \\          % $m_z$（绕$z$轴的力矩）
0 \\          % $f_x$（沿$x$轴的力）
0 \\          % $f_y$（沿$y$轴的力）
\lambda_3     % $f_z$（沿$z$轴的力）
\end{bmatrix}
\end{align}

\textbf{物理意义}：
\begin{itemize}
\item $\lambda_1$：绕$x$轴的约束力矩（防止倾斜）
\item $\lambda_2$：绕$y$轴的约束力矩（防止倾斜）
\item $\lambda_3$：沿$z$轴的约束力（防止穿透黑板）
\end{itemize}

\section{约束力的计算}

\subsection{从动力学方程求解}

\textbf{约束动力学}：
\begin{equation}
F = \Lambda(\theta)\dot{V} + \eta(\theta, V) + A^T(\theta)\lambda
\end{equation}

\textbf{约束的时间导数}：
\begin{equation}
A(\theta)\dot{V} + \dot{A}(\theta)V = 0
\end{equation}

\textbf{求解$\lambda$}：

从约束动力学求解$\dot{V}$：
\begin{equation}
\dot{V} = \Lambda^{-1}(\theta)[F - \eta(\theta, V) - A^T(\theta)\lambda]
\end{equation}

代入约束的时间导数：
\begin{align}
A(\theta)\Lambda^{-1}(\theta)[F - \eta(\theta, V) - A^T(\theta)\lambda] + \dot{A}(\theta)V &= 0 \\
A(\theta)\Lambda^{-1}(\theta)A^T(\theta)\lambda &= A(\theta)\Lambda^{-1}(\theta)[F - \eta(\theta, V)] + \dot{A}(\theta)V
\end{align}

求解$\lambda$：
\begin{equation}
\lambda = [A(\theta)\Lambda^{-1}(\theta)A^T(\theta)]^{-1}\left\{A(\theta)\Lambda^{-1}(\theta)[F - \eta(\theta, V)] + \dot{A}(\theta)V\right\}
\end{equation}

\textbf{约束力}：
\begin{equation}
F_{\text{constraint}} = A^T(\theta)\lambda
\end{equation}

\subsection{矩阵$A\Lambda^{-1}A^T$的性质}

\textbf{矩阵维度}：
\begin{itemize}
\item $A(\theta) \in \mathbb{R}^{k \times 6}$
\item $\Lambda^{-1}(\theta) \in \mathbb{R}^{6 \times 6}$
\item $A^T(\theta) \in \mathbb{R}^{6 \times k}$
\item $A(\theta)\Lambda^{-1}(\theta)A^T(\theta) \in \mathbb{R}^{k \times k}$
\end{itemize}

\textbf{可逆性}：
\begin{itemize}
\item 如果约束是独立的，$A(\theta)$的秩为$k$
\item 如果$\Lambda(\theta)$是正定的，$\Lambda^{-1}(\theta)$也是正定的
\item 那么$A(\theta)\Lambda^{-1}(\theta)A^T(\theta)$是$k \times k$正定矩阵，可逆
\end{itemize}

\section{约束力的物理意义}

\subsection{约束力是什么？}

\textbf{定义}：
\begin{itemize}
\item 约束力是环境施加到机器人上的力
\item 这个力确保机器人满足约束$A(\theta)V = 0$
\item 约束力与约束方向相关
\end{itemize}

\textbf{方向}：
\begin{itemize}
\item 约束力$F_{\text{constraint}}$位于$A^T(\theta)$的列空间中
\item 这是约束子空间的正交补空间
\item 约束力不会影响自由运动方向
\end{itemize}

\textbf{大小}：
\begin{itemize}
\item 约束力的大小由拉格朗日乘数$\lambda$决定
\item $\lambda$的大小取决于：
\begin{itemize}
\item 外部力$F$
\item 科里奥利项$\eta(\theta, V)$
\item 加速度$\dot{V}$
\end{itemize}
\end{itemize}

\subsection{约束力与约束的关系}

\textbf{约束}：$A(\theta)V = 0$

\textbf{约束力}：$F_{\text{constraint}} = A^T(\theta)\lambda$

\textbf{关系}：
\begin{itemize}
\item 约束定义了速度必须满足的条件
\item 约束力确保速度满足约束
\item 约束力的方向与约束方向正交（在约束子空间的正交补空间中）
\end{itemize}

\section{详细例子}

\subsection{例子1：简单约束（$k = 1$）}

\textbf{约束}：$v_z = 0$（不能沿$z$轴运动）

\textbf{约束矩阵}：
\begin{equation}
A(\theta) = \begin{bmatrix} 0 & 0 & 0 & 0 & 0 & 1 \end{bmatrix}
\end{equation}

\textbf{转置矩阵}：
\begin{equation}
A^T(\theta) = \begin{bmatrix} 0 \\ 0 \\ 0 \\ 0 \\ 0 \\ 1 \end{bmatrix}
\end{equation}

\textbf{拉格朗日乘数}：
\begin{equation}
\lambda = \lambda_1 \in \mathbb{R}
\end{equation}

\textbf{约束力}：
\begin{equation}
F_{\text{constraint}} = A^T(\theta)\lambda = \begin{bmatrix} 0 \\ 0 \\ 0 \\ 0 \\ 0 \\ \lambda_1 \end{bmatrix}
\end{equation}

\textbf{物理意义}：
\begin{itemize}
\item 约束力只在$z$方向（$f_z = \lambda_1$）
\item 其他方向的力为零
\item 这确保了$v_z = 0$的约束
\end{itemize}

\subsection{例子2：多个约束（$k = 3$）}

\textbf{约束}：
\begin{align}
\omega_x &= 0 \\
\omega_y &= 0 \\
v_z &= 0
\end{align}

\textbf{约束矩阵}：
\begin{equation}
A(\theta) = \begin{bmatrix}
1 & 0 & 0 & 0 & 0 & 0 \\
0 & 1 & 0 & 0 & 0 & 0 \\
0 & 0 & 0 & 0 & 0 & 1
\end{bmatrix}
\end{equation}

\textbf{转置矩阵}：
\begin{equation}
A^T(\theta) = \begin{bmatrix}
1 & 0 & 0 \\
0 & 1 & 0 \\
0 & 0 & 0 \\
0 & 0 & 0 \\
0 & 0 & 0 \\
0 & 0 & 1
\end{bmatrix}
\end{equation}

\textbf{拉格朗日乘数}：
\begin{equation}
\lambda = \begin{bmatrix} \lambda_1 \\ \lambda_2 \\ \lambda_3 \end{bmatrix}
\end{equation}

\textbf{约束力}：
\begin{equation}
F_{\text{constraint}} = A^T(\theta)\lambda = \begin{bmatrix}
\lambda_1 \\  % $m_x$
\lambda_2 \\  % $m_y$
0 \\          % $m_z$
0 \\          % $f_x$
0 \\          % $f_y$
\lambda_3     % $f_z$
\end{bmatrix}
\end{equation}

\textbf{物理意义}：
\begin{itemize}
\item $\lambda_1$：绕$x$轴的约束力矩（防止$\omega_x \neq 0$）
\item $\lambda_2$：绕$y$轴的约束力矩（防止$\omega_y \neq 0$）
\item $\lambda_3$：沿$z$轴的约束力（防止$v_z \neq 0$）
\end{itemize}

\section{总结}

\subsection{矩阵构型总结}

\begin{enumerate}
\item \textbf{约束矩阵}：$A(\theta) \in \mathbb{R}^{k \times 6}$
\begin{itemize}
\item $k$行：$k$个约束
\item $6$列：6维任务空间（twist）
\end{itemize}

\item \textbf{转置矩阵}：$A^T(\theta) \in \mathbb{R}^{6 \times k}$
\begin{itemize}
\item $6$行：6维任务空间（wrench）
\item $k$列：$k$个约束方向
\end{itemize}

\item \textbf{拉格朗日乘数}：$\lambda \in \mathbb{R}^k$
\begin{itemize}
\item $k$个标量，每个对应一个约束
\end{itemize}

\item \textbf{约束力}：$F_{\text{constraint}} = A^T(\theta)\lambda \in \mathbb{R}^6$
\begin{itemize}
\item 6维wrench（力和力矩）
\item 位于$A^T(\theta)$的列空间中
\end{itemize}
\end{enumerate}

\subsection{关键公式}

\begin{enumerate}
\item \textbf{约束力}：
\begin{equation}
F_{\text{constraint}} = A^T(\theta)\lambda
\end{equation}

\item \textbf{拉格朗日乘数}：
\begin{equation}
\lambda = [A(\theta)\Lambda^{-1}(\theta)A^T(\theta)]^{-1}\left\{A(\theta)\Lambda^{-1}(\theta)[F - \eta(\theta, V)] + \dot{A}(\theta)V\right\}
\end{equation}

\item \textbf{约束动力学}：
\begin{equation}
F = \Lambda(\theta)\dot{V} + \eta(\theta, V) + A^T(\theta)\lambda
\end{equation}
\end{enumerate}

\subsection{物理意义}

\begin{itemize}
\item \textbf{约束力}：环境施加到机器人上的力，确保约束满足
\item \textbf{方向}：约束力位于约束子空间的正交补空间中
\item \textbf{大小}：由拉格朗日乘数$\lambda$决定，取决于外部力、科里奥利项和加速度
\item \textbf{作用}：约束力不会影响自由运动方向，只影响约束方向
\end{itemize}

\end{document}

